% ==================================================
% PREÁMBULO DEL LIBRO
% Aquí van paquetes, colores, comandos y entornos.
% Normalmente NO hace falta tocar este archivo para escribir.
% ==================================================

% Idioma y codificación
\usepackage[utf8]{inputenc}
\usepackage[T1]{fontenc}
\usepackage[spanish,es-nodecimaldot]{babel}

% Tipografía
\usepackage{newtxtext}
\usepackage{newtxmath}
\usepackage{microtype}

% Página
\usepackage[
    a4paper,
    inner=2.8cm,
    outer=2.3cm,
    top=2.5cm,
    bottom=2.5cm,
    headheight=15pt
]{geometry}

% Matemáticas
\usepackage{amsmath,amssymb,amsthm,mathtools}
\usepackage{bm}

% Gráficos y tablas
\usepackage{graphicx}
\usepackage{booktabs}
\usepackage{array}

% Color y cajas
\usepackage[dvipsnames,svgnames]{xcolor}
\usepackage[most]{tcolorbox}
\tcbuselibrary{theorems,breakable,skins}
\usetikzlibrary{arrows.meta}

% Listas
\usepackage{enumitem}

% Cabeceras
\usepackage{fancyhdr}
\usepackage{emptypage}
\usepackage{titlesec}
\usepackage{titletoc}
\usepackage{makeidx}
\makeindex

% ==================================================
% COLORES
% ==================================================

\definecolor{maincolor}{HTML}{5B42F5}
\definecolor{darkcolor}{HTML}{29243D}
\definecolor{lightcolor}{HTML}{F4F2FF}
\definecolor{graycolor}{HTML}{666666}
\definecolor{linecolor}{HTML}{DCD9E8}

% Enlaces (hyperref se carga después de definir los colores)
\usepackage[
    colorlinks=true,
    linkcolor=maincolor,
    citecolor=maincolor,
    urlcolor=maincolor
]{hyperref}

% ==================================================
% FORMATO GENERAL
% ==================================================

\setlength{\parindent}{0pt}
\setlength{\parskip}{5pt}

\setcounter{secnumdepth}{3}
\setcounter{tocdepth}{2}

\setlist[itemize]{leftmargin=1.5em, itemsep=2pt, topsep=4pt}
\setlist[enumerate]{leftmargin=1.8em, itemsep=2pt, topsep=4pt}

% ==================================================
% CAPÍTULOS Y SECCIONES
% ==================================================

\linespread{1.06}

% (Los \par van dentro de macros porque titlesec no admite \par en sus argumentos)
\newcommand{\capfilete}{\textcolor{maincolor}{\rule{2.2cm}{1.2pt}}\par\vspace{1.2em}}
\newcommand{\capetiqueta}{%
  \textcolor{maincolor}{\scshape\large\chaptertitlename\ \thechapter}\par\vspace{0.5em}\capfilete}

% Capítulo numerado: etiqueta pequeña, filete corto y título grande
\titleformat{\chapter}[display]
{\normalfont\raggedright}
{\capetiqueta}
{0pt}
{\bfseries\fontsize{30}{34}\selectfont\color{darkcolor}}

% Capítulo sin numerar (Prefacio, Notación...)
\titleformat{name=\chapter,numberless}[display]
{\normalfont\raggedright}
{}{0pt}
{\capfilete\bfseries\fontsize{30}{34}\selectfont\color{darkcolor}}

\titlespacing*{\chapter}{0pt}{20pt}{34pt}

\titleformat{\section}
{\normalfont\Large\bfseries\color{darkcolor}}
{\textcolor{maincolor}{\thesection}}{0.8em}{}

\titleformat{\subsection}
{\normalfont\large\bfseries\color{darkcolor}}
{\textcolor{maincolor}{\thesubsection}}{0.8em}{}

% ==================================================
% ÍNDICE GENERAL
% ==================================================

\titlecontents{chapter}[2.5em]
{\addvspace{12pt}\bfseries\color{darkcolor}}
{\color{maincolor}\contentslabel{2.5em}}
{\hspace*{-2.5em}}
{\hfill\contentspage}

\titlecontents{section}[5.5em]
{}
{\contentslabel{3em}}
{\hspace*{-3em}}
{\titlerule*[0.6em]{.}\contentspage}

% ==================================================
% COMANDOS PERSONALIZADOS
% (añade aquí los tuyos; así los cambias en un solo sitio)
% ==================================================

\newcommand{\R}{\mathbb{R}}
\newcommand{\N}{\mathbb{N}}
\newcommand{\dd}{\mathrm{d}}
\newcommand{\e}{\mathrm{e}}
\newcommand{\dif}{\mathop{}\!\mathrm{d}}
\newcommand{\der}{\frac{\mathrm{d}}{\mathrm{d}t}}
\newcommand{\pd}{\partial}
\newcommand{\avg}[1]{\left\langle #1\right\rangle}

% Página preliminar suelta (cubierta, portada...) sin páginas en blanco intermedias
\newenvironment{pagina}{\clearpage\thispagestyle{empty}}{\clearpage}

% Separador ornamental dentro de un capítulo
\newcommand{\ornamento}{%
  \begin{center}\textcolor{maincolor}{\small$\bullet\ \ \bullet\ \ \bullet$}\end{center}}

% Epígrafe al inicio de un capítulo: \epigrafe{Texto}{Autor}
\newcommand{\epigrafe}[2]{%
  \begin{flushright}\begin{minipage}{0.7\textwidth}\raggedleft
  \small\itshape #1\par\vspace{3pt}
  {\upshape\textcolor{maincolor}{--- #2}}
  \end{minipage}\end{flushright}\vspace{1.5em}}

% ==================================================
% ENTORNOS (cajas con numeración automática)
% ==================================================
%
% Uso:
%   \begin{definicion}{Título}{etiqueta} ... \end{definicion}
%   -> se referencia con \ref{def:etiqueta}
%
% Todos comparten el mismo contador dentro de cada capítulo:
% Definición 1.1, Teorema 1.2, Ejemplo 1.3, ...

% Estilo común de las cajas con título
\tcbset{
    cajabase/.style={
        enhanced, breakable, arc=0pt,
        left=9pt, right=9pt, top=8pt, bottom=8pt,
        before skip=10pt, after skip=10pt,
        fonttitle=\bfseries, coltitle=darkcolor,
        attach boxed title to top left={yshift*=-\tcboxedtitleheight/2},
    }
}

% Definición
\newtcbtheorem[number within=chapter]{definicion}{Definición}{
    enhanced, breakable,
    colback=lightcolor, colframe=maincolor, colbacktitle=lightcolor,
    coltitle=darkcolor, fonttitle=\bfseries,
    boxrule=0pt, borderline west={2pt}{0pt}{maincolor},
    arc=0pt, left=9pt, right=9pt, top=8pt, bottom=8pt,
    before skip=10pt, after skip=10pt,
    detach title, before upper={\tcbtitle.\ \ },
}{def}

% Teorema
\newtcbtheorem[use counter from=definicion]{teorema}{Teorema}{
    enhanced, breakable,
    colback=white, colframe=linecolor, boxrule=0.5pt,
    borderline west={2pt}{0pt}{maincolor},
    arc=0pt, left=9pt, right=9pt, top=8pt, bottom=8pt,
    before skip=10pt, after skip=10pt,
    fonttitle=\bfseries, coltitle=darkcolor, colbacktitle=white,
    detach title, before upper={\tcbtitle.\ \ },
}{thm}

% Proposición
\newtcbtheorem[use counter from=definicion]{proposicion}{Proposición}{
    enhanced, breakable,
    colback=white, colframe=linecolor, boxrule=0.5pt,
    borderline west={2pt}{0pt}{maincolor},
    arc=0pt, left=9pt, right=9pt, top=8pt, bottom=8pt,
    before skip=10pt, after skip=10pt,
    fonttitle=\bfseries, coltitle=darkcolor, colbacktitle=white,
    detach title, before upper={\tcbtitle.\ \ },
}{prop}

% Lema
\newtcbtheorem[use counter from=definicion]{lema}{Lema}{
    enhanced, breakable,
    colback=white, colframe=linecolor, boxrule=0.5pt,
    borderline west={2pt}{0pt}{maincolor},
    arc=0pt, left=9pt, right=9pt, top=8pt, bottom=8pt,
    before skip=10pt, after skip=10pt,
    fonttitle=\bfseries, coltitle=darkcolor, colbacktitle=white,
    detach title, before upper={\tcbtitle.\ \ },
}{lema}

% Corolario
\newtcbtheorem[use counter from=definicion]{corolario}{Corolario}{
    enhanced, breakable,
    colback=white, colframe=linecolor, boxrule=0.5pt,
    borderline west={2pt}{0pt}{maincolor},
    arc=0pt, left=9pt, right=9pt, top=8pt, bottom=8pt,
    before skip=10pt, after skip=10pt,
    fonttitle=\bfseries, coltitle=darkcolor, colbacktitle=white,
    detach title, before upper={\tcbtitle.\ \ },
}{cor}

% Ejemplo
\newtcbtheorem[use counter from=definicion]{ejemplo}{Ejemplo}{
    enhanced, breakable,
    colback=white, colframe=linecolor, boxrule=0.5pt, arc=2pt,
    left=9pt, right=9pt, top=8pt, bottom=8pt,
    before skip=10pt, after skip=10pt,
    fonttitle=\bfseries, coltitle=darkcolor, colbacktitle=white,
    detach title, before upper={\tcbtitle.\ \ },
}{ej}

% Demostración (sin numeración; añade el cuadradito final sola)
\newtcolorbox{demostracion}[1][Demostración]{
    enhanced, breakable,
    colback=white, colframe=white, boxrule=0pt,
    borderline west={1pt}{0pt}{graycolor},
    arc=0pt, left=12pt, right=9pt, top=5pt, bottom=5pt,
    before skip=8pt, after skip=8pt,
    before upper={\textbf{#1.}\ },
    after upper={\hfill$\square$},
}

% Resultado importante (fondo de color)
\newtcolorbox{resultado}[1][]{
    enhanced, breakable,
    colback=maincolor, colframe=maincolor, coltext=white,
    boxrule=0pt, arc=2pt,
    left=10pt, right=10pt, top=9pt, bottom=9pt,
    before skip=12pt, after skip=12pt,
    #1
}

% Observación (sin numeración)
\newtcolorbox{observacion}[1][Observación]{
    enhanced, breakable,
    colback=gray!6, colframe=gray!6, boxrule=0pt,
    borderline west={2pt}{0pt}{graycolor},
    arc=0pt, left=9pt, right=9pt, top=7pt, bottom=7pt,
    before skip=8pt, after skip=8pt,
    before upper={\textbf{#1.}\ },
}

% ==================================================
% CABECERAS Y PIES
% ==================================================

\pagestyle{fancy}
\fancyhf{}

\renewcommand{\chaptermark}[1]{\markboth{#1}{}}
\renewcommand{\sectionmark}[1]{\markright{\thesection\ \ #1}}

\newcommand{\cabsep}{\hspace{1.2em}{\color{maincolor}\vrule width 0.5pt height 0.9em depth 0.25em}\hspace{1.2em}}

\fancyhead[LE]{\small\thepage\cabsep\scshape\leftmark}
\fancyhead[RO]{\small\scshape\rightmark\cabsep\thepage}

\renewcommand{\headrulewidth}{0.4pt}
\renewcommand{\headrule}{{\color{linecolor}%
  \hrule width\headwidth height\headrulewidth \vskip-\headrulewidth}}
\renewcommand{\footrulewidth}{0pt}

% Primera página de capítulo: solo número de página abajo
\fancypagestyle{plain}{%
  \fancyhf{}%
  \fancyfoot[C]{\small\thepage}%
  \renewcommand{\headrulewidth}{0pt}%
}

% Metadatos del PDF
\hypersetup{
  pdftitle={Mecánica de Distribuciones},
  pdfauthor={Mario Sultan Romero}
}
